\documentclass{article}
\usepackage{graphicx}
\usepackage{wrapfig}

\begin{document}

\section{Wrapped figure alongside text}

\begin{wrapfigure}{r}{0.35\textwidth}
  \centering
  \includegraphics[width=0.3\textwidth]{diagram}
  \caption{An illustrative figure flowed alongside the prose.}
  \label{fig:wrap}
\end{wrapfigure}

The wrapfig package lets a figure be flowed alongside the body text rather than taking a full horizontal slice of the page. The syntax \texttt{wrapfigure\{r\}\{width\}} places the figure on the right with the specified width.

This paragraph demonstrates the wrapping: the text flows around the figure, saving vertical space on the page. Useful for small illustrations that relate directly to a nearby paragraph.

Keep in mind that wrapfig has quirks with line-breaking and should not span across multiple text blocks. See Figure~\ref{fig:wrap} for reference.

After the wrapped region, paragraphs resume full-width.

\end{document}
